\section{Future work}

To address the limitations identified in this study and build upon our findings, three primary avenues for future research are proposed:

First, future work will focus on integrating Large Language Models (LLMs) and Large Vision-Language Models (LVLMs) to transition static recommenders into proactive, conversational AI styling agents. By wrapping graph collaborative filtering backends inside conversational LLM agents, the system can engage in multi-round natural language styling dialogues, as exemplified by unified vision-language fashion assistants such as FashionM3 \parencite{pang2025fashionm3multimodalmultitaskmultiround}. This integration will enable the recommender to explain the rationale behind its outfit recommendations and dynamically adapt suggestions based on conversational feedback, directly addressing the challenges and opportunities of agentic personalized fashion recommendation in the generative AI era \parencite{deldjoo2025agenticpersonalizedfashionrecommendation}.

Second, we propose incorporating trend-aware temporal dynamics into the graph propagation process. Rather than caching static offline visual embeddings, future models can integrate time-varying graph neural networks or temporal attention layers. This will allow the system to model real-time fashion trend evolution, seasonal styling transitions, and sequential user interest drift, ensuring that recommendations remain seasonally relevant and aligned with current market styles.

Finally, to address the computational overhead and scale the framework to massive e-commerce datasets, future architectures will explore scalable and decoupled GNN frameworks. Techniques such as sub-graph sampling, graph partitioning, or decoupling graph propagation from representation learning will be evaluated. This will enable the multimodal recommender engine to operate efficiently on millions of active users and items, maintaining low latency and high scalability without GPU memory exhaustion.
